%=====================================================================
\section{Final --- 2025, Segundo Cuatrimestre (1)}
%=====================================================================

%=====================================================================
\subsection{Ejercicio 1 --- Pulseras Bluetooth (AES-CTR, modo degradado)}
%=====================================================================

\begin{consigna}
Eran pulseras bluetooth para abrir puertas que cada dos segundos mandaban su
$(ID, CTR, C)$ donde $CTR$ es un contador, $C = \mathrm{AES\text{-}CTR}_k(CTR \,\|\,
\mathit{bit\_wearing})$. $\mathit{bit\_wearing}$ es un bit que indica si es o no una
persona (si la pulsera está puesta).

El servidor, al recibir $(ID, CTR, C)$, descifra $C$ con la $K$ asignada a ese $ID$
y compara si el $CTR$ descifrado es el mismo que el que venía en claro; si lo es, le
da el acceso.

Si el servidor se queda sin conexión, se entra en un \textbf{estado degradado} donde
la puerta deja pasar a cualquier $ID$ que haya tenido acceso recientemente, sin
hacer ninguna otra validación.

\begin{enumerate}
  \item[a)] ¿Qué validaciones podría hacer el servidor para evitar comportamientos
  extraños como que se abran dos puertas distantes casi al mismo tiempo, o que se
  intenten abrir puertas fuera de un horario normal?
  \item[b)] Si alguien captura los mensajes, explicar cómo podría utilizarlos o
  modificarlos para tener acceso. ¿Cómo lo solucionaría?
  \item[c)] Si consiguen el estado degradado, ¿qué podría hacerse para que no puedan
  acceder repitiendo mensajes?
\end{enumerate}
\end{consigna}

\paragraph{Temas.}
Primer Parte, Clase 2 --- Cifrado (\S Modo CTR: keystream $E_k(\text{nonce}\|i)$,
nunca usa $D_k$, paralelizable, CPA-Secure solo si no se repite $(k,\text{nonce})$);
Primer Parte, Clase 3 --- MACs y Cifrado Autenticado (\S Maleabilidad de cifrados de
flujo/CTR; \S CCA y por qué el cifrado solo no alcanza para integridad; \S MAC:
infalsificabilidad, procesar todo el mensaje); Segunda Parte, Clase 2 ---
Autenticación (\S OTP/HOTP: contador sincronizado, códigos de un solo uso,
throttling); Segunda Parte, Clase 5 --- Seguridad de la Empresa (\S IDS/IPS:
detección por anomalías; \S SIEM: correlación de eventos y alertas).

\paragraph{Qué necesitás saber.}
\begin{itemize}
  \item \textbf{CTR es un cifrado de flujo, no autenticado.} $C_i = E_k(\text{nonce}
  \|i)\oplus M_i$; el descifrado es $M_i = E_k(\text{nonce}\|i)\oplus C_i$ (nunca se
  usa $D_k$). CTR da \emph{confidencialidad} (si no se repite el par clave/nonce),
  pero \textbf{no} da integridad ni autenticación por sí solo.
  \item \textbf{Maleabilidad de cifrados de flujo (Clase 3).} Si
  $\mathrm{Enc}_k(m) = G(k)\oplus m = c$, entonces $c' = c\oplus x \Rightarrow
  \mathrm{Dec}_k(c') = m\oplus x$: un atacante que conoce (o adivina) un bit de
  plano en una posición conocida puede flippear ese bit del cifrado, y al descifrar
  cambia exactamente ese bit del plano, \textbf{sin conocer la clave}. Es el mismo
  mecanismo que el ataque al sueldo en RR.HH.\ del material.
  \item \textbf{Ningún criptosistema de los vistos es CCA-Secure}: el cifrado está
  pensado para confidencialidad, no para resistir manipulación del texto cifrado.
  Por eso la integridad requiere una primitiva aparte: el \textbf{MAC}. Un MAC
  robusto siempre procesa la \emph{totalidad} del mensaje (si no, lo que queda
  afuera se puede alterar sin invalidar la etiqueta).
  \item \textbf{OTP/HOTP (Clase 2, Segunda Parte):} $A$ y $B$ comparten una clave
  $K$ y un contador $C$; el código depende de $(K,C)$ y se recomienda
  \emph{throttling} y (para HOTP) una ventana de \emph{look-ahead} para tolerar
  desincronización del contador.
  \item \textbf{IDS/IPS y SIEM (Clase 5, Segunda Parte):} un IDS/IPS detecta por
  \textbf{anomalías} (además de firmas) y decide bloquear o alertar; un SIEM
  \textbf{correlaciona eventos} de distintas fuentes (aquí, aperturas de puertas
  distintas) para detectar patrones anómalos y generar alertas.
\end{itemize}

\paragraph{Notas y conexiones.}
El enunciado describe un esquema de \emph{confidencialidad sin integridad}: se cifra
$(CTR\|\mathit{bit\_wearing})$ con AES-CTR, pero el mensaje completo $(ID,CTR,C)$ no
lleva ningún MAC que lo ate. Esto es exactamente el escenario que motiva la Clase 3
(``el cifrado por sí solo no protege contra modificaciones''): acá, además, el
$CTR$ viaja \textbf{en claro} junto al cifrado, lo que facilita un replay directo. La
verificación descripta en el enunciado (``compara si el CTR es el mismo que el que
venía en $C$'') es débil porque no dice qué pasa si el $CTR$ ya fue usado antes
(no hay noción de \emph{frescura/monotonicidad}), y no ata $ID$ al contenido cifrado
con una clave de integridad.

\paragraph{Resolución.}

\textbf{a) Validaciones de contexto (correlación entre eventos, no solo por mensaje).}
Estas validaciones no dependen de la criptografía del mensaje individual sino de
\textbf{cruzar eventos en el servidor} (tipo SIEM: correlación de eventos + reglas
de anomalía, IDS/IPS basado en anomalías):
\begin{itemize}
  \item \textbf{Verificación de coherencia física (velocidad de desplazamiento):} el
  servidor conoce la posición de cada puerta. Si el mismo $ID$ abrió una puerta y
  poco después (menos tiempo del físicamente necesario para recorrer la distancia)
  intenta abrir otra puerta lejana, la solicitud es \textbf{físicamente imposible}
  y debe rechazarse y generar una alerta (correlación de eventos por $ID$ a través
  del tiempo y el espacio).
  \item \textbf{Ventanas horarias (ACL temporal):} asociar a cada $ID$ (o a cada
  puerta) un rango horario habilitado en una base de datos; si llega una trama fuera
  de ese rango, denegar el acceso y generar una alerta de seguridad, en lugar de
  autorizar silenciosamente.
  \item En ambos casos el mecanismo es \textbf{detectivo/correlativo} (como un
  IDS/SIEM): no reemplaza la validación criptográfica del mensaje individual, la
  complementa detectando patrones anómalos entre mensajes.
\end{itemize}

\textbf{b) Captura de mensajes: uso/modificación y solución.}
Como $C$ se obtiene con AES-CTR (un cifrado de flujo), aplica directamente la
\textbf{maleabilidad de los cifrados de flujo} vista en la Clase 3:
\begin{itemize}
  \item \textbf{Replay puro:} el atacante captura un mensaje válido $(ID, CTR, C)$ y
  lo reenvía más tarde. El enunciado solo dice que el servidor ``compara si el CTR
  es el mismo que el que viene en $C$'' --- \emph{no} dice que el servidor lleve un
  registro del último $CTR$ usado por ese $ID$. Si no lo lleva, el mensaje capturado
  \textbf{se vuelve a aceptar tal cual} (replay exitoso).
  \item \textbf{Modificación por maleabilidad (si el atacante conoce/adivina el bit
  de \emph{wearing}):} dado que CTR es un XOR con el keystream
  ($C = E_k(\text{nonce}\|i) \oplus (CTR\|\mathit{bit\_wearing})$), sin conocer la
  clave $k$ el atacante puede flippear el bit del cifrado que corresponde a la
  posición del bit $\mathit{bit\_wearing}$ dentro de $C$: eso cambia exactamente ese
  bit al descifrar (igual que el ataque al sueldo: $c'' = c\oplus x$ cambia el bit
  correspondiente del plano), \textbf{sin necesitar la clave}. Así podría forzar
  a que el servidor interprete que la pulsera está puesta (o no) según convenga
  para pasar alguna validación adicional que dependa de ese bit.
  \item \textbf{Solución:} agregar \textbf{integridad} con un MAC sobre el mensaje
  completo $(ID\|CTR\|\mathit{bit\_wearing})$, calculado con una clave de MAC
  (idealmente distinta de la de cifrado) y verificado antes de aceptar el mensaje.
  Un MAC que cubre la totalidad del contenido evita tanto la modificación por
  maleabilidad (cualquier bit flippeado invalida la etiqueta) como facilita atar el
  $ID$ al contenido cifrado. Además, el servidor debe llevar un \textbf{registro del
  último $CTR$ aceptado por $ID$} y solo aceptar $CTR_{\text{nuevo}} > CTR_{\text{último}}$
  (monotonicidad / frescura), lo que bloquea el replay puro.
\end{itemize}

\textbf{c) Estado degradado: evitar replay sin conexión.}
En el modo degradado el servidor no puede verificar nada dinámicamente (no tiene
conexión para consultar estado), así que la defensa tiene que apoyarse en algo que
\textbf{se consuma} localmente, análogo a la idea de \textbf{OTP/HOTP} (Clase 2,
Segunda Parte): claves de un solo uso derivadas de un contador compartido.
\begin{itemize}
  \item Precargar en la memoria de la puerta (mientras hay conexión) una
  \textbf{lista de códigos de un solo uso} (equivalente a un HOTP: valores derivados
  de $(K, C)$ con $C$ incremental), vinculados a cada pulsera/$ID$.
  \item Cuando no hay conexión, la pulsera envía el \textbf{siguiente código no usado}
  de esa lista y la puerta lo valida \textbf{localmente} y lo \textbf{marca como
  consumido} (borrándolo o avanzando un puntero), de forma que no pueda reutilizarse:
  esto es exactamente la garantía de un esquema de un solo uso (OTP), que rompe el
  replay porque cada código sirve una única vez.
  \item Al recuperar la conexión, el servidor debe re-sincronizar el estado
  (equivalente al \emph{look-ahead} de HOTP) y reponer/renovar la lista de códigos
  consumidos localmente.
\end{itemize}

%=====================================================================
\subsection{Ejercicio 2 --- Protocolo de autenticación Cliente--Servidor ($H(K_c\|r)$)}
%=====================================================================

\begin{consigna}
Protocolo de autenticación entre un Cliente $C$ y un Servidor $S$ que ambos tienen
una $K_c$ compartida. La función $H$ es un hash sin colisiones.
\begin{align*}
C \to S &: ID\\
S \to C &: r \quad (\text{número aleatorio})\\
C \to S &: H(K_c \| r)
\end{align*}
El servidor calcula el hash usando la $K_c$ que tiene guardada y el $r$. Si da lo
mismo que recibió, autentica a $C$.

Preguntas:
\begin{enumerate}
  \item[1.] ¿Por qué el servidor manda ese $r$? ¿Qué pasaría si el cliente solo
  mandase $H(K_c)$ en vez de que se use el $r$?
  \item[2.] ¿Disminuiría la seguridad si se usara un contador incremental en vez
  del $r$ aleatorio?
  \item[3.] Plantear cómo se podría hacer para que el cliente autentique también al
  servidor, sin introducir nuevas claves.
\end{enumerate}
\end{consigna}

\paragraph{Temas.}
Segunda Parte, Clase 2 --- Autenticación (\S Challenge-Response y EKE: el desafío
$r$ como nonce, ``el secreto nunca viaja por la red''); Primer Parte, Clase 5 ---
Protocolos (\S Ataque por reflexión y ``romper la simetría'' con nonces/identidad
por sentido; \S Handshakes con nonces de ambos lados, ej.\ TLS \emph{Hello} con
$r_1, r_2$).

\paragraph{Qué necesitás saber.}
\begin{itemize}
  \item \textbf{Challenge-response (Clase 2, Segunda Parte):} el servidor manda un
  \textbf{challenge único} $r$ (nonce); el cliente responde $f(k,r)$ (acá,
  $H(K_c\|r)$) probando que conoce el secreto $K_c$ \textbf{sin enviarlo}. El
  diagrama del material remarca explícitamente: ``el nonce evita reutilizar
  respuestas capturadas''.
  \item Si no hubiera $r$ (o el cliente mandara siempre $H(K_c)$, un valor fijo), la
  respuesta sería \textbf{siempre la misma}: un atacante que la capture una vez
  puede \textbf{reenviarla} (replay) y autenticarse sin conocer $K_c$.
  \item Un \textbf{contador incremental} en vez de un $r$ aleatorio también da
  frescura (cada valor se usa una sola vez si el servidor lo verifica creciente),
  pero es \textbf{predecible}: un atacante que observa la secuencia puede anticipar
  el próximo valor. En el material esto es exactamente el punto que distingue HOTP
  (contador, requiere sincronización, predecible en su progresión) de un valor
  \textbf{aleatorio} imprevisible.
  \item \textbf{Autenticación mutua sin nuevas claves:} el material resuelve el
  problema simétrico en los \textbf{handshakes} (Clase 5, Primer Parte) generando
  un desafío \textbf{en cada sentido} con la misma clave compartida, y remarca
  (a propósito del ataque por reflexión) que hay que \textbf{romper la simetría}
  entre ambos desafíos (por ejemplo, atando la identidad/dirección a cada
  desafío) para que una sesión paralela no pueda usarse para reflejar la
  respuesta.
\end{itemize}

\paragraph{Notas y conexiones.}
Este protocolo es el patrón \textbf{challenge-response} de la Clase 2 (Segunda
Parte) aplicado con hash en lugar de una función $f(k,r)$ genérica. El ataque por
\textbf{reflexión} de la Clase 5 (Primer Parte) es la advertencia clave para la
pregunta 3: si simplemente se agrega un segundo desafío simétrico (el cliente
manda un $r_c$ y espera $H(K_c\|r_c)$ del servidor, de la misma forma en que el
servidor se lo pide al cliente), un atacante podría abrir una \textbf{sesión
paralela} y reflejar el desafío recibido en una sesión para responder en la otra,
si no se distingue el rol/dirección de cada desafío.

\paragraph{Resolución.}

\textbf{1) Por qué se manda $r$; qué pasa con $H(K_c)$ solo.}
El servidor manda $r$ para dar \textbf{frescura} a la respuesta: como $r$ es
distinto en cada intento de autenticación, la respuesta $H(K_c\|r)$ también lo es,
así que \textbf{capturar una respuesta pasada no sirve para autenticarse de
nuevo} (no hay dos ejecuciones del protocolo con la misma respuesta válida). Si
el cliente mandara directamente $H(K_c)$ (sin usar $r$), la respuesta sería
\textbf{siempre el mismo valor fijo} en cada autenticación; un atacante que la
capture una sola vez (escuchando el canal) puede luego \textbf{reenviarla
(replay)} y el servidor la aceptaría como si fuera el cliente legítimo, sin
conocer nunca $K_c$. Es exactamente el rol del nonce en challenge-response: ``el
secreto nunca viaja por la red'' y \textbf{cada respuesta es única}, ligada a un
desafío que no se repite.

\textbf{2) Contador incremental en vez de $r$ aleatorio.}
Sí \textbf{disminuye la seguridad} respecto de un nonce verdaderamente aleatorio,
aunque igual aportaría algo de frescura si el servidor exige que sea creciente
(evitaría el replay literal del mismo mensaje). El problema es la
\textbf{predictibilidad}: si el contador es incremental y el atacante observa la
secuencia (por ejemplo, escuchando varias rondas), puede \textbf{anticipar el
valor del próximo contador} y, si además logra capturar u obtener de otra forma el
valor de $H(K_c\|r)$ para un $r$ futuro conocido (o construir un ataque que
aproveche esa previsibilidad, según qué tan expuesto esté el mecanismo), reduce el
espacio de incertidumbre del protocolo. En el material esto es análogo a la
comparación entre HOTP (contador, requiere sincronización explícita y su
progresión es previsible) y un desafío aleatorio (TOTP/nonce, imprevisible salvo
que se conozca la clave). Un valor \textbf{aleatorio e imprevisible} es preferible
porque no permite a un atacante preparar de antemano una respuesta para un
desafío futuro.

\textbf{3) Autenticación mutua sin introducir nuevas claves.}
Usando la misma $K_c$ compartida, se agrega un desafío \textbf{en el otro
sentido}, rompiendo la simetría para evitar un ataque por reflexión (Clase 5,
Primer Parte):
\begin{align*}
C \to S &: r_C,\ ID\\
S \to C &: r_S,\ H(K_c \| r_C)\\
C \to S &: H(K_c \| r_S)
\end{align*}
El servidor prueba que conoce $K_c$ respondiendo al desafío $r_C$ del cliente
(en el mismo mensaje en que manda su propio desafío $r_S$); el cliente, a su vez,
responde al desafío $r_S$. Cada parte verifica el hash correspondiente a \emph{su
propio} desafío. Es importante que ambos desafíos \textbf{no sean intercambiables
uno por el otro} sin ambigüedad (en el material, la defensa contra la reflexión es
distinguir la dirección/rol de cada nonce, por ejemplo incluyendo dentro del hash
la identidad o el rol del que responde, del estilo $H(K_c\|r_S\|\text{``cliente''})$
y $H(K_c\|r_C\|\text{``servidor''})$) para que una sesión paralela abierta por un
atacante no pueda usarse para que el servidor legítimo ``firme'' un desafío que
luego se reutilice contra el cliente. No se introduce ninguna clave nueva: se
reutiliza $K_c$ y $H$ en ambos sentidos.

%=====================================================================
\subsection{Ejercicio 3 --- CasaIA (IoT hogareño): hipótesis de falla}
%=====================================================================

\begin{consigna}
Una página web y aplicación mobile conectadas con una API REST. Se llamaba CasaIA,
te dejaba controlar los dispositivos de tu casa a través de wifi, y podrías
configurar funciones por horario y por ubicación geográfica del teléfono. Los
dispositivos están conectados con un websocket en TCP al servidor, que es el que
les puede mandar comandos (abrir/cerrar, etc.).

Los dispositivos se vinculan escaneando un QR que viene en ellos.

Se puede dar acceso temporal a través de la página a otra persona: el usuario pone
el mail de la persona y a esta le llega un mail con el cual, al abrirlo, consigue
acceso temporal. También se genera un token de sesión al lograrse el acceso.

Se pide: como siempre, dar \textbf{4 hipótesis de falla} y su prueba.
\end{consigna}

\paragraph{Temas.}
Segunda Parte, Clase 6 --- Pentesting (\S Metodología de hipótesis de falla, Paso 2:
áreas candidatas de hipótesis; \S la prueba debe ser repetible y concluyente);
Segunda Parte, Clase 3 --- Principios de Diseño y Vulnerabilidades (\S CSRF; \S XSS;
\S nonce/token anti-replay); Segunda Parte, Clase 2 --- Autenticación (\S OTP/HOTP,
contador vs.\ aleatorio; \S challenge-response); glosario (\S OAuth2/tokens,
\S nonce).

\paragraph{Qué necesitás saber.}
\begin{itemize}
  \item \textbf{Metodología de hipótesis de falla (Clase 6):} el paso 2 consiste en
  plantear posibles vulnerabilidades \textbf{examinando implementaciones y
  mecanismos} (¿está bien implementado el control?, ¿el ambiente introduce
  errores?, ¿es realmente seguro?) y \textbf{comparando con sistemas parecidos}
  (sistemas similares tienen problemas similares). El paso 3 (prueba) exige que la
  verificación de la hipótesis sea \textbf{repetible y concluyente}.
  \item \textbf{Nonce / frescura (Clase 3):} la defensa real contra CSRF es un
  \textbf{token/nonce único} que dé frescura y evite replay; lo mismo aplica a
  cualquier enlace o token de un solo uso (como el mail de acceso temporal de
  CasaIA).
  \item \textbf{XSS (Clase 3):} inyectar código que termine ejecutado por el
  browser de la víctima permite \textbf{robar la sesión} (aprovechando la confianza
  del sitio en su sesión/cookie).
  \item \textbf{Tokens de sesión (glosario):} un token de sesión, si no tiene
  expiración adecuada, alcance limitado, o viaja sin protección, puede ser robado o
  reutilizado para suplantar al usuario.
  \item \textbf{OTP/HOTP (Clase 2):} un código de acceso temporal enviado por mail
  que actúa como un ``one-time'' debería, como en HOTP/TOTP, tener \textbf{throttling}
  y no ser reutilizable ni de vida indefinida.
\end{itemize}

\paragraph{Nota sobre alcance (regla de oro).}
El enunciado pide una \textbf{metodología puntual} (``4 hipótesis de falla, una
prueba cada una''); en el material de Clase 6 (Pentesting) la metodología está
descripta en \textbf{5 pasos} (Recolección $\to$ Hipótesis $\to$ Prueba $\to$
Generalización $\to$ Eliminación) y el paso de \textbf{Hipótesis} es explícitamente
donde se listan posibles vulnerabilidades (no hay, en las clases provistas, una
regla de puntaje del tipo ``4 hipótesis, 1 punto cada una''; eso parece ser una
convención de la consigna del parcial, no un contenido de la clase). Lo que sí está
en el material es \emph{cómo} se plantea y se prueba cada hipótesis (paso 2 y 3 de
la metodología), que es lo que se aplica a continuación sobre CasaIA.

\paragraph{Resolución.}
Aplicando la metodología de hipótesis de falla (examinar implementación/mecanismos
+ comparar con sistemas parecidos) a los cuatro componentes distintivos de CasaIA
(QR de vinculación, websocket de comandos, acceso temporal por mail, token de
sesión):

\begin{enumerate}
  \item \textbf{Hipótesis --- vinculación por QR sin verificación de posesión.} Si
  el QR del dispositivo solo contiene un identificador estático (sin un secreto o
  desafío que se consuma una única vez), cualquiera que \textbf{fotografíe o copie}
  el QR (por ejemplo, publicado sin querer en una foto, o visto por un instalador)
  podría vincular ese dispositivo a su propia cuenta. \emph{Prueba (repetible y
  concluyente, Clase 6):} tomar una foto del QR de un dispositivo ya vinculado y,
  desde otra cuenta, intentar vincularlo nuevamente; si el sistema lo permite (o no
  invalida la vinculación anterior), la hipótesis queda confirmada.
  \item \textbf{Hipótesis --- falta de autenticación/autorización en los mensajes
  del websocket.} Si el canal websocket entre servidor y dispositivo no reautentica
  cada comando (solo se autenticó la conexión TCP inicial) ni valida que el comando
  provenga de un usuario autorizado sobre \emph{ese} dispositivo puntual, un
  atacante que logre interceptar o inyectar tráfico en ese canal podría mandar
  comandos (abrir/cerrar) sin ser el dueño. \emph{Prueba:} con acceso a la red del
  dispositivo (caja gris), capturar el tráfico del websocket y reenviar
  (\emph{replay}) un comando ``abrir'' capturado, o inyectar un comando propio, y
  verificar si el servidor lo ejecuta sin re-validar autorización.
  \item \textbf{Hipótesis --- el enlace de acceso temporal por mail no tiene
  frescura/expiración adecuada (análogo a CSRF/replay, Clase 3).} Si el link que
  llega por mail para dar acceso temporal es un token que no expira, no es de un
  solo uso, o es predecible/corto, cualquiera que lo intercepte (compartido reenvío
  de mail, mail comprometido, o el link cacheado en el navegador) podría usarlo
  para obtener acceso indefinidamente, igual que un token anti-CSRF débil no da
  frescura real. \emph{Prueba:} generar un acceso temporal, usarlo una vez, y
  volver a usar el mismo enlace pasado un tiempo prudencial (o tras revocar el
  acceso desde la app); si sigue funcionando, la hipótesis queda confirmada.
  \item \textbf{Hipótesis --- robo/reutilización del token de sesión (estilo
  XSS/robo de sesión, Clase 3).} Si el token de sesión generado al lograrse el
  acceso (temporal o del dueño) no tiene alcance ni expiración acotados, o si la
  interfaz web es vulnerable a inyección de código (XSS) que permita leer ese
  token, un atacante podría robarlo y \textbf{suplantar la sesión} sin conocer
  ninguna contraseña. \emph{Prueba:} intentar inyectar un script en algún campo de
  la aplicación web (por ejemplo, el nombre del dispositivo o de la persona
  invitada) y verificar si se ejecuta al ser renderizado por otro usuario
  (confirmando que el token/cookie de sesión de ese usuario quedaría expuesto);
  alternativamente, verificar si el token de sesión sigue siendo válido después de
  que expiró el acceso temporal que lo originó.
\end{enumerate}

